\documentclass[12pt,a4paper]{article}

%--------------- Packages -----------------
\usepackage{graphicx}
\usepackage{amsmath}
\usepackage{booktabs}
\usepackage{array}
\usepackage{geometry}
\usepackage{hyperref}
\usepackage[table]{xcolor}
\usepackage{enumitem}

\geometry{margin=1in}

\begin{document}
	
	%%%%%%%%%%%%%%%%%%%%%%%%%%%%%
	% Header
	%%%%%%%%%%%%%%%%%%%%%%%%%%%%%
	
	\begin{center}
		{\Large \textbf{MALAVIYA NATIONAL INSTITUTE OF TECHNOLOGY, JAIPUR}}\\[0.2cm]
		{\large Department of Computer Science and Engineering}\\[0.3cm]
		{\Large \textbf{B.Tech III Semester Examination}}\\[0.2cm]
		{\large Course: Introduction to LaTeX}\\[0.5cm]
	\end{center}
	
	\noindent
	\textbf{Time: 2 Hours} \hfill \textbf{Maximum Marks: 30}
	
	\vspace{0.5cm}
	
	\noindent
	\textbf{Student Name:} \rule{6cm}{0.4pt}
	
	\vspace{0.3cm}
	
	\noindent
	\textbf{Roll Number:} \rule{5cm}{0.4pt}
	
	\hrule
	\vspace{0.4cm}
	
	%%%%%%%%%%%%%%%%%%%%%%%%%%%%%
	% Instructions
	%%%%%%%%%%%%%%%%%%%%%%%%%%%%%
	
	\section*{Instructions}
	
	\begin{enumerate}
		\item Attempt all questions.
		\item Write all answers using LaTeX syntax.
		\item Figures and tables should have captions.
		\item Marks are indicated against each question.
	\end{enumerate}
	
	%%%%%%%%%%%%%%%%%%%%%%%%%%%%%
	% Section A
	%%%%%%%%%%%%%%%%%%%%%%%%%%%%%
	
	\section*{Section A (10 Marks)}
	
	\begin{enumerate}[label=\textbf{Q\arabic*.}]
		
		\item Explain the following LaTeX commands with examples. \hfill (4 Marks)
		
		\begin{itemize}
			\item \texttt{\\section}
			\item \texttt{\\subsection}
			\item \texttt{\\textbf}
			\item \texttt{\\emph}
		\end{itemize}
		
		\item Write LaTeX code to create the following equation. \hfill (3 Marks)
		
		\[
		E = mc^2
		\]
		
		\item Create an ordered and unordered list in LaTeX. \hfill (3 Marks)
		
	\end{enumerate}
	
	%%%%%%%%%%%%%%%%%%%%%%%%%%%%%
	% Section B
	%%%%%%%%%%%%%%%%%%%%%%%%%%%%%
	
	\section*{Section B (12 Marks)}
	
	\begin{enumerate}[label=\textbf{Q\arabic*.},start=4]
		
		\item Create the following table using LaTeX. \hfill (6 Marks)
		
		\begin{table}[h]
			\centering
			\caption{Student Marks}
			\begin{tabular}{|c|c|c|}
				\hline
				\rowcolor{gray!20}
				Name & Subject & Marks\\
				\hline
				Ketan & DSA & 95\\
				Rahul & OOP & 90\\
				Aman & DBMS & 88\\
				\hline
			\end{tabular}
		\end{table}
		
		\item Insert an image and give it a caption. \hfill (3 Marks)
		
		\begin{figure}[h]
			\centering
			\includegraphics[width=0.45\textwidth]{1.jpeg}
			\caption{Sample Figure}
			\label{fig:sample}
		\end{figure}
		
		\item Write the difference between \texttt{table} and \texttt{tabular} environments. \hfill (3 Marks)
		
	\end{enumerate}
	
	%%%%%%%%%%%%%%%%%%%%%%%%%%%%%
	% Section C
	%%%%%%%%%%%%%%%%%%%%%%%%%%%%%
	
	\section*{Section C (8 Marks)}
	
	\begin{enumerate}[label=\textbf{Q\arabic*.},start=7]
		
		\item Create a bibliography using \texttt{thebibliography}. Cite one reference in the text. \hfill (4 Marks)
		
		Example citation: Machine learning is widely used \cite{mlbook}.
		
		\begin{thebibliography}{99}
			
			\bibitem{mlbook}
			Tom Mitchell,
			\textit{Machine Learning},
			McGraw Hill, 1997.
			
		\end{thebibliography}
		
		\item Write LaTeX code for the following mathematical expression. \hfill (4 Marks)
		
		\[
		\int_0^1 x^2\,dx = \frac{1}{3}
		\]
		
	\end{enumerate}
	
	%%%%%%%%%%%%%%%%%%%%%%%%%%%%%
	% Footer
	%%%%%%%%%%%%%%%%%%%%%%%%%%%%%
	
	\vfill
	\hrule
	
	\begin{center}
		\textbf{--- End of Question Paper ---}
	\end{center}
	
\end{document}